% Generated draft: requires mathematical and source review.
\documentclass[11pt]{article}
\usepackage[margin=25mm]{geometry}
\usepackage{fontspec}
\setmainfont{Latin Modern Roman}
\usepackage{amsmath,amssymb,mathtools,unicode-math}
\setmathfont{Latin Modern Math}
\usepackage{enumitem,needspace,xurl,hyperref}
\setlength{\parindent}{0pt}
\setlength{\parskip}{5pt}
\setlength{\emergencystretch}{3em}
\newcommand{\N}{\mathbb{N}}
\newcommand{\Z}{\mathbb{Z}}
\newcommand{\Q}{\mathbb{Q}}
\newcommand{\R}{\mathbb{R}}
\newcommand{\C}{\mathbb{C}}
\newcommand{\F}{\mathbb{F}}
\newcommand{\A}{\mathbb{A}}
\newcommand{\PP}{\mathbb{P}}
\newcommand{\E}{\mathbb{E}}
\newcommand{\eps}{\varepsilon}
\newcommand{\dd}{\,\mathrm{d}}
\begin{document}
\section*{UnsolvedMath ID: 700777}
\textbf{Source problem code:} catalogue entry
\subsection{Definitions and mathematical statement}
Let \(p\) be a prime, and let \(F\) be a free pro-\(p\) group with free pro-\(p\) basis \(X\). No finiteness assumption on \(X\) is supplied. For \(t\in F\), write
\[
t^F=\overline{\langle f^{-1}tf\mid f\in F\rangle}
\]
for its closed normal closure. Fix \(r\in F\), put \(R=r^F\), and let \(q:F\to F/R\) be the quotient homomorphism. An element \(s\in F\) is \emph{associated to \(r\)} if \(s^F=R\).

\textbf{(a)} Suppose that
\[
\forall s\in F,\qquad s^F=R\ \Longrightarrow\ \neg\exists u\in F\;(s=u^p).
\]
Must \(F/R\) be torsion-free? Explicitly, must
\[
\forall g\in F/R\;\forall n\in\Z_{\geq 1},\qquad g^n=1\ \Longrightarrow\ g=1
\]
hold?

\textbf{(b)} For \(Y\subseteq X\), let \(F_Y=\overline{\langle Y\rangle}\leq F\). The natural interpretation of an element depending on \(Y\) is that it belongs to \(F_Y\). Under this interpretation, define
\[
\mathcal S_r=\{Y\subseteq X\mid \exists s\in F_Y\text{ with }s^F=R\}.
\]
The source states that an element associated to \(r\) depends on a minimal subset \(X'\) of \(X\). For such an \(X'\), and each \(x\in X'\), must the images \(q(y)\), indexed by \(y\in X\setminus\{x\}\), freely generate a free pro-\(p\) group? Equivalently, is the restriction
\[
q\big|_{F_{X\setminus\{x\}}}:F_{X\setminus\{x\}}\longrightarrow F/R
\]
injective, or, equivalently, does
\[
R\cap F_{X\setminus\{x\}}=\{1\}
\]
hold?

The source does not define dependence or specify the meaning of minimality. In particular, minimality under inclusion in \(\mathcal S_r\), minimum cardinality among members of \(\mathcal S_r\), and minimal support for a particular associated element are not distinguished. Thus part (b) retains an unresolved qualification: the displayed question expresses its requested conclusion, but a fully quantified formulation requires clarification of the choice of \(X'\). The existence of such a minimal subset is stated in the source and is not established here. Part (b) is presented separately; the source does not explicitly say that the hypothesis of part (a) carries over.
\subsection{Short English statement}
For a free pro-\(p\) group, part (a) asks whether the quotient by one closed normal relation is torsion-free when no element generating the same closed normal subgroup is a \(p\)th power. Part (b) asks whether deleting a basis element occurring in a suitably minimal associated relation leaves freely generating images in the quotient. The meaning of minimality in part (b) requires clarification. The supplied registry labels the problem open; this draft does not independently verify that status.
\subsection{Sources}
\begin{itemize}
\item Local UnsolvedMath registry, permanent ID \(700777\), exact source code \texttt{catalogue entry}, title \emph{Kourovka Notebook Problem 11.85}. Supplied upstream revision: \texttt{d371ee917ac9647fbb06cd6fc529deec6f3dda82}. The record was created and updated on \texttt{2026-10-16T14:00:00Z}; these are registry timestamps, not literature-verification dates.
\item \emph{The Kourovka Notebook (Unsolved Problems in Group Theory)}, original problem \(11.85\), 11th issue (1990), as identified in the supplied record through the 21st issue edition (2026), arXiv identifier \texttt{1401.0300v48}. The supplied collection description names E. I. Khukhro and V. D. Mazurov as editors. This bibliographic identification is supplied evidence and has not been independently checked.
\item Proposed in 1990. The proposer field is supplied as \texttt{N. S. Romanovski\textbackslash{}u\{\textbackslash{}i}; its final TeX diacritic is incomplete, so the spelling is retained without repair.
\end{itemize}
\end{document}
